% Standalone problem document, generated from the adjacent README.md.
% Compile with XeLaTeX. Mathematical content is maintained in Markdown.
\documentclass[11pt,a4paper]{article}
\usepackage[fontsize=10.5pt]{scrextend}
\usepackage[margin=22mm,headheight=15pt,headsep=9mm,footskip=11mm]{geometry}
\usepackage{fontspec}
\setmainfont{texgyrepagella}[Extension=.otf,UprightFont=*-regular,BoldFont=*-bold,ItalicFont=*-italic,BoldItalicFont=*-bolditalic]
\setsansfont{texgyreheros}[Extension=.otf,UprightFont=*-regular,BoldFont=*-bold,ItalicFont=*-italic,BoldItalicFont=*-bolditalic]
\setmonofont{texgyrecursor}[Extension=.otf,UprightFont=*-regular,BoldFont=*-bold,ItalicFont=*-italic,BoldItalicFont=*-bolditalic,Scale=MatchLowercase]
\usepackage{amsmath,amssymb}
\usepackage{unicode-math}
\setmathfont{texgyrepagella-math.otf}
\usepackage{xcolor,microtype,enumitem,needspace,fancyhdr}
\usepackage{bookmark}
\definecolor{ink}{HTML}{17344B}
\definecolor{accent}{HTML}{197D80}
\definecolor{muted}{HTML}{596773}
\hypersetup{colorlinks=true,linkcolor=accent,urlcolor=accent,pdftitle={TR-14 - Comon's
exact-rank conjecture for Hankel
tensors},pdfauthor={Open Problems in Numerical Linear Algebra}}
\urlstyle{same}
\setlength{\parindent}{0pt}
\setlength{\parskip}{5pt plus 1pt minus 1pt}
\linespread{1.04}
\setlength{\emergencystretch}{3em}
\widowpenalty=10000
\clubpenalty=10000
\displaywidowpenalty=10000
\allowdisplaybreaks[1]
\setlist{nosep,leftmargin=1.5em}
\providecommand{\tightlist}{\setlength{\itemsep}{0pt}\setlength{\parskip}{0pt}}
\providecommand{\pandocbounded}[1]{#1}
\pagestyle{fancy}
\fancyhf{}
\fancyhead[L]{\sffamily\footnotesize\color{muted}OPEN PROBLEMS IN NUMERICAL LINEAR ALGEBRA}
\fancyhead[R]{\sffamily\bfseries\footnotesize\color{accent}TR-14}
\fancyfoot[L]{\parbox[b]{0.9\textwidth}{\sffamily\fontsize{8}{10}\selectfont\color{muted}Editorial ratings; a bounded search cannot certify that a problem remains unsolved.\\Literature check: 2026-09-08}}
\fancyfoot[R]{\sffamily\footnotesize\color{muted}\thepage}
\renewcommand{\headrulewidth}{0.3pt}
\renewcommand{\headrule}{\hbox to\headwidth{\color{accent}\leaders\hrule height \headrulewidth\hfill}}
\makeatletter
\renewcommand{\section}{\Needspace{5\baselineskip}\@startsection{section}{1}{0pt}{12pt plus 2pt minus 2pt}{4pt}{\sffamily\bfseries\large\color{ink}}}
\renewcommand{\subsection}{\Needspace{4\baselineskip}\@startsection{subsection}{2}{0pt}{9pt plus 1pt minus 1pt}{3pt}{\sffamily\bfseries\normalsize\color{ink}}}
\makeatother
\setcounter{secnumdepth}{0}
\begin{document}
{\sffamily\small\color{muted}Tensor computations\par}
\vspace{3pt}
{\raggedright\hyphenpenalty=10000\exhyphenpenalty=10000\sffamily\bfseries\fontsize{21}{24}\selectfont\color{ink}Comon's
exact-rank conjecture for Hankel tensors\par}
\vspace{7pt}
{\sffamily\small\textbf{Difficulty:} challenging\quad\textbf{Importance:} interesting
to the community\par}
\vspace{4pt}
{\color{accent}\hrule height 0.8pt}
\vspace{6pt}
\subsection{Statement}\label{statement}

For every \(m\ge3\), \(n\ge2\), and \(h\in\mathbb C^{m(n-1)+1}\), form
the symmetric tensor \[
H_{i_1\ldots i_m}=h_{i_1+\cdots+i_m-m},\qquad 1\le i_j\le n.
\] Is \(R(H)=R_{\rm sym}(H)\) always true? Here \(R(H)\) is the least
integer \(r\ge0\) admitting \[
H=\sum_{j=1}^r u_{j,1}\otimes\cdots\otimes u_{j,m},\quad u_{j,k}\in\mathbb C^n,
\] whereas \(R_{\rm sym}(H)\) is the least \(r\) admitting
\(H=\sum_{j=1}^r c_jv_j^{\otimes m}\), with \(c_j\in\mathbb C\) and
\(v_j\in\mathbb C^n\). Empty sums represent zero. No genericity or
Vandermonde restriction is imposed.

\subsection{Relevance}\label{relevance}

This would justify symmetry-preserving decomposition lengths for every
Hankel tensor, including exceptional data. Such structure appears in
harmonic retrieval and structured tensor factorization.

\subsection{References}\label{references}

\begin{enumerate}
\def\labelenumi{\arabic{enumi}.}
\tightlist
\item
  J. Nie and K. Ye, \emph{Hankel Tensor Decompositions and Ranks}, SIAM
  J. Matrix Anal. Appl. 40 (2019).
  \href{https://epubs.siam.org/doi/10.1137/18M1168285}{DOI};
  \href{https://arxiv.org/pdf/1706.03631}{primary preprint}, §6,
  Conjecture 6.3, p.16.
\item
  L. Qi, \emph{Hankel Tensors: Associated Hankel Matrices and
  Vandermonde Decomposition}, 2014.
  \href{https://arxiv.org/pdf/1310.5470}{Primary preprint}, §1, equation
  (1), and §4, Theorem 3, for Hankel structure and decomposition.
\end{enumerate}

\subsection{Status check --- 2026-09-08}\label{status-check-2026-09-08}

Nie and Ye formulate this restricted conjecture after explicitly
acknowledging Shitov's counterexample to the unrestricted
symmetric-tensor conjecture. Searches
\textit{"Hankel" "Conjecture 6.3"},
\textit{"Hankel" "Comon" "counterexample"}, and
\textit{"Hankel tensors" ranks conjectures solved} located no
solution. No later primary reaffirmation was located; the dated
conclusion is bounded by this search.
\end{document}
